%==============================================================================
%  MetaLab Suite --- Problem #1 (beginner edition)
%  Wind Farm Layout Optimization: a first metaheuristics exercise
%
%  Compile with:  pdflatex wind_farm_layout_optimization.tex   (twice for the ToC)
%==============================================================================
\documentclass[11pt,a4paper]{article}

\usepackage[T1]{fontenc}
\usepackage[utf8]{inputenc}
\usepackage{lmodern}
\usepackage[margin=2.5cm]{geometry}
\usepackage{amsmath,amssymb,mathtools}
\usepackage{booktabs}
\usepackage{algorithm}
\usepackage{algpseudocode}
\usepackage{xcolor}
\usepackage{microtype}
\usepackage[hidelinks]{hyperref}

\newcommand{\R}{\mathbb{R}}
\DeclareMathOperator*{\argmax}{arg\,max}

\title{%
  \textbf{Wind Farm Layout Optimization}\\[2mm]
  \large A First Metaheuristics Exercise}
\author{%
  MetaLab Suite\\[1mm]
  \small Problem \#1 --- beginner edition}
\date{\today}

%==============================================================================
\begin{document}
\maketitle

\begin{abstract}
\noindent
This is the friendly, first-days-of-the-course version of the Wind Farm Layout
Optimization (WFLO) problem. We strip the problem down to a small grid, two
wind directions, a single wake rule and one objective, so that the only thing
left to do is the thing we actually want to teach: \emph{how to design a
metaheuristic}. By the end you will have implemented your first hill climber and
compared it against random search and the exact optimum. The full, continuous
version of the problem is described in the advanced companion document.
\end{abstract}

\tableofcontents

%==============================================================================
\section{The problem in words}
%==============================================================================

A wind farm is a small \textbf{grid of cells}. We own $N$ turbines and must
decide \textbf{which $N$ cells} to occupy. The catch is the \emph{wake effect}:
a turbine slows the air behind it, so any turbine standing downwind of it
receives less wind, and therefore produces less power. Placing turbines far
apart reduces this interference, but we cannot make the farm infinitely large.

To make the problem concrete and small enough to reason about, we use a
``toy'' model with exactly two wind directions:

\begin{itemize}
  \item the wind blows \textbf{east} (left $\rightarrow$ right) with
        probability $p_{\text{east}}$;
  \item the wind blows \textbf{south} (top $\rightarrow$ bottom) with
        probability $p_{\text{south}} = 1 - p_{\text{east}}$.
\end{itemize}

The two directions are what make the problem interesting. With only the east
wind, the best plan would be to line every turbine up in the \emph{leftmost
column} (nothing is upwind of anything). But the south wind would then knock
those very turbines down: standing in a single column, they shelter one another
vertically. The optimal layout must \textbf{stagger} the turbines so that the
wake losses are shared fairly between the two directions.

%==============================================================================
\section{A simple model}
%==============================================================================

\subsection{Domain and decision variables}

The farm is an $m\times m$ grid of cells. Number the cells
$0,1,\dots,m^{2}-1$; cell $c$ sits at row and column
\begin{equation}
  (r, y) = (c \div m,\ c \bmod m).
\end{equation}
A \textbf{layout} is a set
\begin{equation}
  X = \{c_{1}, c_{2}, \dots, c_{N}\} \subset \{0, 1, \dots, m^{2}-1\},
  \qquad \lvert X \rvert = N,
\end{equation}
i.e.\ a choice of $N$ distinct cells. This set of cells is the only decision we
make.

\subsection{The wake rule}

Wind direction determines the direction of the wakes, and therefore which
turbines are ``upwind'' of which. We use one fixed loss factor
$\lambda_{\text{wake}} = \texttt{wake\_loss} \in [0,1]$ and the simple rule:

\begin{quote}
A turbine loses a fraction $\lambda_{\text{wake}}$ of its wind speed
\emph{for every turbine that is upwind of it on its own line}.
\end{quote}

``Upwind on its own line'' means, concretely,
\begin{itemize}
  \item \textbf{east wind:} other turbines in the \emph{same row} with a
        \emph{smaller column};
  \item \textbf{south wind:} other turbines in the \emph{same column} with a
        \emph{smaller row}.
\end{itemize}

Let $k_{i}^{d}$ be the number of turbines upwind of turbine $i$ for direction
$d \in \{\text{east},\ \text{south}\}$. The speed turbine $i$ sees is
\begin{equation}
  u_{i}^{d} = u_{\infty}\,\bigl(1-\lambda_{\text{wake}}\bigr)^{\,k_{i}^{d}},
  \label{eq:speed}
\end{equation}
where $u_{\infty} = \texttt{free\_speed}$ is the free-stream wind speed. Being
second on a line costs one factor $(1-\lambda_{\text{wake}})$, being third costs
two, and so on.

\subsection{Turbine power}

Each turbine converts wind speed into electrical power through the cubic power
curve
\begin{equation}
  P(u) = P_{r}\left(\min\!\Bigl(1,\ \tfrac{u}{u_{r}}\Bigr)\right)^{3},
  \label{eq:power}
\end{equation}
where $u_{r} = \texttt{rated\_speed}$ and $P_{r} = \texttt{rated\_power}$ are
the rated speed and rated power. Below the rated speed the power grows with the
cube of the wind ($P \propto u^{3}$); above it the turbine is capped.

\subsection{Objective: expected total power}

For one direction, the farm produces the sum of its turbines' power,
\begin{equation}
  P^{d}(X) = \sum_{i=1}^{N} P\!\left(u_{i}^{d}(X)\right),
\end{equation}
and the quantity we want to \textbf{maximise} is the expected power over the two
directions:
\begin{equation}
  \boxed{\;
  P_{\text{total}}(X)
  = p_{\text{east}}\, P^{\text{east}}(X)
  + p_{\text{south}}\, P^{\text{south}}(X)
  = \sum_{d} p_{d} \sum_{i=1}^{N} P\!\left(u_{i}^{d}(X)\right). \;}
  \label{eq:objective}
\end{equation}
The complete problem is
\begin{equation}
  \max_{X} \ P_{\text{total}}(X)
  \quad\text{subject to}\quad
  X \subset \{0,\dots,m^{2}-1\},\ \lvert X\rvert = N .
  \label{eq:problem}
\end{equation}

\begin{quote}
\textbf{Constraint (only one).} The $N$ turbines occupy \emph{distinct} cells.
Everything else is objective.
\end{quote}

%==============================================================================
\section{Why we need a metaheuristic}
%==============================================================================

\subsection{The search space is exponential}

The number of possible layouts is the number of ways to choose $N$ cells out of
$m^{2}$:
\begin{equation}
  \lvert \mathcal{X} \rvert = \binom{m^{2}}{N}.
  \label{eq:space}
\end{equation}
For the default instance ($m=4$, $N=6$) this is only
$\binom{16}{6} = 8008$ --- small enough that we \emph{can} check every layout by
brute force, which is wonderful for grading. But the space explodes: going to
$m=8$, $N=16$ gives $\binom{64}{16} \approx 4.9\times10^{14}$ layouts, far
beyond enumeration.

\subsection{The objective is a nasty, non-smooth landscape}

Even though the problem looks tame, the objective \eqref{eq:objective} has two
features that defeat simple, greedy reasoning:
\begin{itemize}
  \item \textbf{coupling.} Whether a turbine suffers a wake depends on
        \emph{which other} turbines were placed in its row and column, so the
        contribution of a turbine cannot be optimised in isolation.
  \item \textbf{multimodality.} Because the two directions pull in different
        ways, good layouts come in competing ``families'' (staggered patterns
        aligned one way or the other), separated by mediocre layouts. There are
        many local optima.
\end{itemize}
Trying to improve a layout one greedy step at a time gets trapped in whichever
local optimum it happens to fall into. This is exactly the situation
\emph{metaheuristics} are built for.

%==============================================================================
\section{Your first metaheuristic: hill climbing}
%==============================================================================

A \emph{trajectory-based} metaheuristic walks a single solution through the
search space. Hill climbing is the simplest one:

\begin{enumerate}
  \item start from a random layout;
  \item make a small random change to get a \textbf{neighbouring} layout;
  \item if the neighbour is better, keep it; otherwise throw it away;
  \item repeat many times, remembering the best layout ever seen.
\end{enumerate}

Two ingredients define any such method: the \textbf{neighbourhood} (what a
``small change'' is) and the \textbf{acceptance rule} (when to keep the change).
For hill climbing the acceptance rule is simply: \emph{keep it if it is
better}. The neighbourhood is the part we leave to you.

\subsection{A neighbourhood for this problem}

A natural move is to take \textbf{one turbine} and move it to a
\textbf{currently empty} cell:
\begin{equation}
  \text{neighbour}(X):\quad
  \text{pick } i \text{ at random},\quad
  \text{set } c_{i} \leftarrow c \text{ with } c \notin X .
  \label{eq:move}
\end{equation}
This always produces a valid layout (still $N$ distinct cells), which is one of
the reasons this encoding is so pleasant.

\begin{algorithm}[H]
\caption{Hill climbing (maximisation)}
\label{alg:hc}
\begin{algorithmic}[1]
\State $X \leftarrow$ \Call{RandomLayout}{$m, N$}
\State $X^{\star} \leftarrow X$ \Comment{best layout so far}
\For{$t = 1$ \textbf{to} $T$}
  \State $X' \leftarrow$ \Call{Neighbour}{$X$} \Comment{the move you implement}
  \If{$P_{\text{total}}(X') > P_{\text{total}}(X)$}
      \State $X \leftarrow X'$
      \If{$P_{\text{total}}(X) > P_{\text{total}}(X^{\star})$}
          \State $X^{\star} \leftarrow X$
      \EndIf
  \EndIf
\EndFor
\State \Return $X^{\star}$
\end{algorithmic}
\end{algorithm}

%==============================================================================
\section{Your task}
%==============================================================================

Everything is provided except the neighbourhood. Open
\texttt{hill\_climbing.py} and implement \texttt{random\_neighbour()}: move one
random turbine to a random empty cell and return the new layout. Then:

\begin{enumerate}
  \item run hill climbing and compare its best power with random search;
  \item compare both with the exhaustive optimum (\texttt{model.brute\_force});
  \item read the suggestions below and try one of them.
\end{enumerate}

\begin{table}[H]
\centering
\caption{Where each piece lives (folder \texttt{starter/}).}
\label{tab:files}
\begin{tabular}{@{}ll@{}}
\toprule
File & Role \\
\midrule
\texttt{model.py}          & given: instance, wake rule, objective, brute force \\
\texttt{generate.py}       & given: build instances (e.g.\ $m=4$, $N=6$) \\
\texttt{random\_search.py} & given: baseline you must beat \\
\texttt{hill\_climbing.py} & \textbf{your task}: implement the neighbourhood \\
\texttt{app.py}            & given: interactive grid visualisation \\
\bottomrule
\end{tabular}
\end{table}

\paragraph{Stretch goal --- simulated annealing.}
Hill climbing only ever accepts improvements, so it stops as soon as no
neighbour is better. \textbf{Simulated annealing} instead accepts a worsening
move with probability $\exp(-\Delta / T)$, where $\Delta>0$ is how much worse
the neighbour is and $T$ is a \emph{temperature} that slowly decreases. Early on
it wanders freely (high $T$); later it behaves like hill climbing (low $T$).
Changing the acceptance rule in Algorithm~\ref{alg:hc} to this rule is a
one-line experiment with a very visible effect on the convergence curve.

%==============================================================================
\section{Further reading}
%==============================================================================

The advanced companion document derives the continuous, physically detailed
version of this problem, including the Jensen wake model, the full power curve,
topographical constraints, and the mathematical justification for using
trajectory metaheuristics. The classical references are \cite{jensen1983} for
the wake model and \cite{mosetti1994,grady2005} for early metaheuristic
approaches to turbine placement. Simulated annealing is due to
\cite{kirkpatrick1983}.

%==============================================================================
\begin{thebibliography}{9}

\bibitem{jensen1983}
N.~O.~Jensen,
``A note on wind generator interaction,''
Ris\o{} National Laboratory, Roskilde, Denmark, Tech.\ Rep.\ Ris\o-M-2411, 1983.

\bibitem{mosetti1994}
G.~Mosetti, C.~Poloni, and B.~Diviacco,
``Optimization of wind turbine positioning in large wind farms by means of a
genetic algorithm,''
\emph{Journal of Wind Engineering and Industrial Aerodynamics},
vol.~51, no.~1, pp.~105--116, 1994.

\bibitem{grady2005}
S.~A.~Grady, M.~Y.~Hussaini, and M.~M.~Abdullah,
``Placement of wind turbines using genetic algorithms,''
\emph{Renewable Energy}, vol.~30, no.~2, pp.~259--270, 2005.

\bibitem{kirkpatrick1983}
S.~Kirkpatrick, C.~D.~Gelatt, and M.~P.~Vecchi,
``Optimization by simulated annealing,''
\emph{Science}, vol.~220, no.~4598, pp.~671--680, 1983.

\end{thebibliography}

\end{document}
